\section{Nhiệm vụ đồ án}\label{sec:nhiem_vu_do_an}
    Đồ án được tổ chức thành 07 chương chính và phần phụ lục, trình bày theo mạch lập luận từ việc nhận diện thách thức thực tế đến đề xuất kiến trúc, hiện thực hóa và kiểm chứng thực nghiệm sản phẩm:

    \begin{itemize}
        \item \textbf{Chương 1: Giới thiệu.} Tập trung phân tích ý tưởng dựa trên các giải pháp mạng máy tính thành công, chỉ rõ những hạn chế của các nền tảng IoT hiện nay về chi phí và tính linh hoạt, từ đó đề xuất giải pháp Gateway IoT dạng module.

        \item \textbf{Chương 2: Lý thuyết.} Trình bày các vấn đề kỹ thuật then chốt cần giải quyết, hệ thống hóa cơ sở lý thuyết về IoT Gateway và các tiêu chuẩn thiết kế phần cứng nâng cao như toàn vẹn tín hiệu (SI), toàn vẹn nguồn (PI).

        \item \textbf{Chương 3: Thiết kế tổng quan.} Đặc tả các yêu cầu sản phẩm và luận giải lựa chọn kiến trúc Dual-MCU ESP32-S3 nhằm tối ưu hóa hiệu năng đa giao thức đồng thời với chi phí thấp.

        \item \textbf{Chương 4: Thiết kế và thực hiện phần cứng.} Mô tả chi tiết quy trình thiết kế mạch nguyên lý và layout PCB cho bo mạch chủ và các module adapter.

        \item \textbf{Chương 5: Thiết kế và thực hiện phần mềm.} Mô tả chi tiết hiện thực hệ thống phần mềm phân lớp, trọng tâm là cơ chế "Software-Defined Module" dựa trên tệp mô tả JSON.

        \item \textbf{Chương 6: Kiểm thử.} Trình bày quy trình kiểm chứng nghiêm ngặt từ mức vật lý bo mạch (Bring-up), đo đạc chất lượng nguồn điện thực tế đến việc kiểm thử chức năng firmware và các chỉ số hiệu suất như thông lượng, độ trễ.

        \item \textbf{Chương 7: Kết quả thực hiện.} Đánh giá kết quả hiện thực hóa sản phẩm thực tế, phân tích hiệu năng vận hành so với lý thuyết và thực hiện đối chiếu chi phí, tính ổn định với các giải pháp thương mại tham chiếu.

        \item \textbf{Chương 8: Kết luận và hướng phát triển.} Tổng kết các mục tiêu đã đạt được, nhận diện các thách thức kỹ thuật trong quá trình thực hiện và đề xuất hướng phát triển sản phẩm tiệm cận các tiêu chuẩn công nghiệp.

        \item \textbf{Phụ lục.} Cung cấp toàn bộ mã nguồn, sơ đồ thiết kế chi tiết, kết quả mô phỏng chuyên sâu và mô tả các cải tiến kỹ thuật quan trọng trong phiên bản Đồ án Tốt nghiệp.
    \end{itemize}

    Trong khuôn khổ đồ án tốt nghiệp, nhóm tập trung vào việc nghiên cứu và hiện thực hóa một mô hình Gateway IoT có tính linh hoạt cao. Tuy nhiên, để đảm bảo tính tập trung và khả năng hoàn thành trong thời gian quy định, đề tài được xác định các giới hạn về phạm vi thực hiện như sau:

    \begin{itemize}
        \item \textbf{Về Kiến trúc Phần cứng:} Đồ án tập trung vào việc thiết kế bo mạch chủ (Baseboard) và các module chuyển đổi (Adapter) dựa trên nền tảng vi điều khiển (MCU) nhằm giảm thiểu chi phí, đáp ứng yêu cầu phân tích và chứng minh tính khả thi của kiến trúc. Do đó, ngoài việc nghiên cứu các kiến trúc Gateway hiệu năng cao dựa trên vi xử lý (MPU) chạy hệ điều hành Linux đầy đủ, việc hiện thực một board mạch như vậy không nằm trong phạm vi của đề tài này.

        \item \textbf{Về Chủng loại module hỗ trợ:} Sản phẩm hiện tại giới hạn hỗ trợ 03 khe cắm vật lý chuẩn hóa. Các giao thức truyền thông được hiện thực và kiểm chứng chỉ bao gồm các chuẩn phổ biến: WAN (Wi-Fi, Ethernet, LTE 4G) và LAN (Zigbee, LoRa, BLE, RS-485). Các giao thức chuyên biệt khác như giao tiếp vệ tinh hay các mạng fieldbus đặc thù chưa được tích hợp trong phiên bản này.

        \item \textbf{Về Phạm vi Xử lý của Firmware:} Hệ thống tập trung giải quyết bài toán vận chuyển dữ liệu (Data Transport), chuyển đổi giao thức và cơ chế cấu hình động qua tệp JSON. Các tính năng xử lý dữ liệu phức tạp tại biên (Edge Computing) như học máy (Machine Learning) hay các thuật toán phân tích dữ liệu chuyên sâu (Data Analytics) nằm ngoài phạm vi thực hiện.

        \item \textbf{Về Quy mô Quản lý Mạng:} Đồ án thực hiện kiểm chứng khả năng vận hành ổn định trong phạm vi từ 20 đến 30 thiết bị đầu cuối kết nối đồng thời. Việc nghiên cứu các giải pháp quản lý mạng quy mô lớn (Enterprise-scale) với hàng ngàn node hoặc quản lý cấu trúc Mesh đa tầng phức tạp nằm ngoài mục tiêu đồ án nhắm tới.

        \item \textbf{Về Tiêu chuẩn Đóng gói:} Sản phẩm dừng lại ở mức độ hoàn thiện nguyên mẫu chức năng có vỏ bảo vệ chống nhiễu. Đồ án không bao gồm việc thực hiện các quy trình kiểm định đạt chuẩn công nghiệp về an toàn điện, tiêu chuẩn viễn thông quốc tế hay các cấp độ bảo vệ môi trường. Các nội dung này sẽ được hiện thực trong phiên bản tiếp theo trong trường hợp có thể đưa sản phẩm ra thị trường.
    \end{itemize}
